\section{Arrays und Standardeingabe}

\subsection{Arrays}
\begin{frame}[fragile]
\frametitle{Arrays}
Was ist ein Array?
\begin{itemize}
    \item speichert eine genaue Anzahl an Variablen desselben Datentypes
    \item hat eine feste Größe (nicht dynamisch erweiterbar)
    \item nicht initialiserte Felder kriegen einen Standardwert oder null zugewiesen
    \item indices für schnellen Zugriff
    \item wird von Java wie ein Objekt behandelt $\rightarrow$ liegt auf dem Heap
\end{itemize}
\begin{lstlisting}
    int[] array1 = new int[5];
    String[] array2 = {"A","R","R","A","Y"};
    System.out.println(array2[0]); // gibt "A" aus
    array1[3] = 42; // array1 = [0,0,0,42,0]
\end{lstlisting}
\end{frame}

\subsection{Mehrdimenionale Arrays}
\begin{frame}[fragile]
\frametitle{Mehrdimenionale Arrays}
Arrays können auch mehrdimensional sein. Das bedeutet, dass wir Arrays von Arrays haben können. Ein zweidimensionales Array könnte beispielsweise so aussehen:
\begin{lstlisting}
    int[][] matrix = {{1, 2},{3, 4},{5, 6}};
    System.out.println(matrix[2][0]); // gibt 5 aus

\end{lstlisting}
\end{frame}